\chapter{2003年河南卷}
\section{单选题}
\begin{problem}
若圆 \(C\) 与圆 \( (x+2)^{2}+(y-1)^{2}=1 \) 关于原点对称，则圆 \(C\) 的方程为
\choices{\( (x-2)^{2}+(y+1)^{2}=1 \)}{\( (x-2)^{2}+(y-1)^{2}=1 \)}{\( (x-1)^{2}+(y+2)^{2}=1 \)}{\( (x+1)^{2}+(y-2)^{2}=1 \)}
\end{problem}

\begin{problem}
抛物线 \(y=ax^{2}\) 的准线方程是 \(y=2\)，则 \(a\) 的值为
\choices{\( \frac{1}{8} \)}{\( -\frac{1}{8} \)}{\(8\)}{\(-8\)}
\end{problem}

\begin{problem}
已知 \( x \in \left(-\frac{\pi}{2}, 0\right) \)， \( \cos x = \frac{4}{5} \)，则 \( \tan 2x = \)
\choices{\( \frac{7}{24} \)}{\( -\frac{7}{24} \)}{\( \frac{24}{7} \)}{\( -\frac{24}{7} \)}
\end{problem}

\begin{problem}
已知四边形 \(ABCD\) 是菱形，点 \(P\) 在对角线 \(AC\) 上（不包括端点 \(A\)，\(C\)），则 \( \overrightarrow{AP} = \)
\choices{\( \lambda(\overrightarrow{AB} + \overrightarrow{AD}) \)（\( \lambda \in (0, 1) \)）}{\( \lambda(\overrightarrow{AB} + \overrightarrow{BC}) \)（\( \lambda \in (0, \frac{\sqrt{2}}{2}) \)）}{\( \lambda(\overrightarrow{AB} - \overrightarrow{AD}) \)（\( \lambda \in (0, 1) \)）}{\( \lambda(\overrightarrow{AB} - \overrightarrow{BC}) \)（\( \lambda \in (0, \frac{\sqrt{2}}{2}) \)）}
\end{problem}

\begin{problem}
设函数 \( f(x)=\begin{cases}
2^{-x}-1, & x\leq 0,\\
\frac{1}{x^{2}}, & x>0.
\end{cases} \)，若 \( f(x_{0})>1 \)，则 \( x_{0} \) 的取值范围是
\choices{\( (-1,1) \)}{\( (-1,+\infty) \)}{\( (-\infty,-2)\cup(0,+\infty) \)}{\( (-\infty,-1)\cup(1,+\infty) \)}
\end{problem}

\begin{problem}
等差数列 \( \{a_{n}\} \) 中，已知 \( a_{1}=\frac{1}{3} \)， \( a_{2}+a_{5}=4 \)， \( a_{n}=33 \)，则 \(n\) 为
\choices{\(48\)}{\(49\)}{\(50\)}{\(51\)}
\end{problem}

\begin{problem}
函数 \( y = \ln \frac{x + 1}{x - 1} \)， \( x \in (1, +\infty) \) 的反函数为
\choices{\( y = \frac{\e^{x} - 1}{\e^{x} + 1} \)（\( x \in (0, +\infty) \)）}{\( y = \frac{\e^{x} + 1}{\e^{x} - 1} \)（\( x \in (0, +\infty) \)）}{\( y = \frac{\e^{x} - 1}{\e^{x} + 1} \)（\( x \in (-\infty, 0) \)）}{\( y = \frac{\e^{x} + 1}{\e^{x} - 1} \)（\( x \in (-\infty, 0) \)）}
\end{problem}

\begin{problem}
棱长为 \(a\) 的正方体中，连结相邻面的中心，以这些线段为棱的八面体的体积为
\choices{\( \frac{a^{3}}{3} \)}{\( \frac{a^{3}}{4} \)}{\( \frac{a^{3}}{6} \)}{\( \frac{a^{3}}{12} \)}
\end{problem}

\begin{problem}
设 \( a > 0 \)， \( f(x) = ax^{2} + bx + c \)，曲线 \( y = f(x) \) 在点 \( P(x_{0}, f(x_{0})) \) 处的切线的倾斜角的取值范围为 \( [0, \frac{\pi}{4}] \)，则 \(P\) 到曲线 \( y = f(x) \) 对称轴距离的取值范围为
\choices{\( \left[0,\frac{1}{a}\right] \)}{\( \left[0,\frac{1}{2a}\right] \)}{\( \left[0,\left|\frac{b}{2a}\right|\right] \)}{\( \left[0,\left|\frac{b-1}{2a}\right|\right] \)}
\end{problem}

\begin{problem}
已知双曲线中心在原点且一个焦点为 \(F(\sqrt{7}, 0)\)，直线 \(y = x - 1\) 与其相交于 \(M\)，\(N\) 两点，\(MN\) 中点的横坐标为 \(-\frac{2}{3}\)，则此双曲线的方程是
\choices{\( \frac{x^{2}}{3}-\frac{y^{2}}{4}=1 \)}{\( \frac{x^{2}}{4}-\frac{y^{2}}{3}=1 \)}{\( \frac{x^{2}}{5}-\frac{y^{2}}{2}=1 \)}{\( \frac{x^{2}}{2}-\frac{y^{2}}{5}=1 \)}
\end{problem}

\begin{problem}
已知长方形的四个顶点 \(A(0, 0)\)， \(B(2, 0)\)， \(C(2, 1)\) 和 \(D(0, 1)\). 一质点从 \(AB\) 的中点 \(P_0\) 沿与 \(AB\) 夹角为 \(\theta\) 的方向射到 \(BC\) 上的点 \(P_1\) 后，依次反射到 \(CD\)， \(DA\) 和 \(AB\) 上的点 \(P_2\)， \(P_3\) 和 \(P_4\)（入射角等于反射角）. 设 \(P_4\) 的坐标为 \( (x_4, 0) \)，若 \( 1 < x_4 < 2 \)，则 \( \tan \theta \) 的取值范围是
\choices{\( \left(\frac{1}{3},1\right) \)}{\( \left(\frac{1}{3},\frac{2}{3}\right) \)}{\( \left(\frac{2}{5},\frac{1}{2}\right) \)}{\( \left(\frac{2}{5},\frac{2}{3}\right) \)}
\end{problem}

\begin{problem}
一个四面体的所有棱长都为 \( \sqrt{2} \)，四个顶点在同一球面上，则此球的表面积为
\choices{\( 3\pi \)}{\( 4\pi \)}{\( 3\sqrt{3}\pi \)}{\( 6\pi \)}
\end{problem}

\section{填空题}
\begin{problem}
\( (x^{2}-\frac{1}{2x})^{9} \) 展开式中 \( x^{9} \) 的系数是\fillinblank{}.
\end{problem}

\begin{problem}
某公司生产三种型号的轿车，产量分别为 \(1200\) 辆，\(6000\) 辆和 \(2000\) 辆. 为检验该公司的产品质量，现用分层抽样的方法抽取 \(46\) 辆进行检验，这三种型号的轿车依次应抽取\fillinblank{}，\fillinblank{}，\fillinblank{}辆.
\end{problem}

\begin{problem}
对于四面体 \(ABCD\)，给出下列四个命题：

\begin{circlelist}
\item 若 \(AB=AC\)，\(BD=CD\)，则 \(BC \perp AD\)；
\item 若 \(AB=CD\)，\(AC=BD\)，则 \(BC \perp AD\)；
\item 若 \(AB \perp AC\)，\(BD \perp CD\)，则 \(BC \perp AD\)；
\item 若 \(AB \perp CD\)，\(BD \perp AC\)，则 \(BC \perp AD\).
\end{circlelist}

其中真命题的序号是\fillinblank{}.（写出所有真命题的序号）
\end{problem}

\begin{problem}
将三种作物种植在如图的 \(5\) 块试验田里，每块种植一种作物且相邻的试验田不能种植同一作物，不同的种植方法共有\fillinblank{}种（以数字作答）.

\begin{center}
\bitmapfigure[max width=0.35\linewidth]{img/2003/henan/q16_fig1.png}
\end{center}
\end{problem}

\section{解答题}
\begin{problem}
已知函数 \( f(x)=2\sin x(\sin x+\cos x) \).

\begin{enumerate}
\item 求函数 \(f(x)\) 的最小正周期和最大值；
\item 在给出的直角坐标系中，画出函数 \(y=f(x)\) 在区间 \(\left[-\frac{\pi}{2}, \frac{\pi}{2}\right]\) 上的图象.

\begin{center}
\bitmapfigure[float=inline,max width=0.60\linewidth]{img/2003/henan/q17_fig1.png}
\end{center}
\end{enumerate}
\end{problem}

\begin{problem}
有三种产品，合格率分别是 \(0.90\)，\(0.95\) 和 \(0.95\)，各抽取一件进行检验.

\begin{enumerate}
\item 求恰有一件不合格的概率；
\item 求至少有两件不合格的概率.（精确到 \(0.001\)）
\end{enumerate}
\end{problem}

\begin{problem}
如图，在直三棱柱 \(ABC-A_{1}B_{1}C_{1}\) 中，底面是等腰直角三角形，\(\angle ACB=90^{\circ}\)，侧棱 \(AA_{1}=2\)，\(D\)，\(E\) 分别是 \(CC_{1}\) 与 \(A_{1}B\) 的中点，点 \(E\) 在平面 \(ABD\) 上的射影是 \(\triangle ABD\) 的重心 \(G\).

\begin{enumerate}
\item 求 \(A_{1}B\) 与平面 \(ABD\) 所成角的大小（结果用反三角函数值表示）；
\item 求点 \(A_{1}\) 到平面 \(AED\) 的距离.

\begin{center}
\bitmapfigure[max width=0.34\linewidth]{img/2003/henan/q19_fig1.png}
\end{center}
\end{enumerate}
\end{problem}

\begin{problem}
已知 \(c > 0\)，设 \(P\)：函数 \(y = c^{x}\) 在 \(\R\) 上单调递减，\(Q\)：不等式 \(x + |x - 2c| > 1\) 的解集为 \(\R\). 如果 \(P\) 和 \(Q\) 有且仅有一个正确，求 \(c\) 的取值范围.
\end{problem}

\begin{problem}
已知常数 \(a > 0\)，向量 \( \bs{c} = (0, a) \)， \( \bs{i} = (1, 0) \)，经过原点 \(O\) 以 \( \bs{c} + \lambda \bs{i} \) 为方向向量的直线与经过定点 \(A(0, a)\) 以 \( \bs{i} - 2\lambda \bs{c} \) 为方向向量的直线相交于点 \(P\)，其中 \(\lambda \in \R\). 试问：是否存在两个定点 \(E\)，\(F\)，使得 \(|PE| + |PF|\) 为定值. 若存在，求出 \(E\)，\(F\) 的坐标；若不存在，说明理由.
\end{problem}

\begin{problem}
已知 \(a > 0\)，\(n\) 为正整数.

\begin{enumerate}
\item 设 \(y=(x-a)^{n}\)，证明 \(y'=n(x-a)^{n-1}\)；
\item 设 \(f_{n}(x)=x^{n}-(x-a)^{n}\)，对任意 \(n \geq a\)，证明 \(f_{n+1}^{\prime}(n+1) > (n+1)f_{n}^{\prime}(n)\).
\end{enumerate}
\end{problem}
